\documentclass[12pt,a4paper]{amsart}
%\input xy
%\xyoption{all}
\usepackage[applemac]{inputenc}

\DeclareMathAlphabet{\mathpzc}{OT1}{pzc}{m}{it}
\usepackage{amsfonts,amsmath,amssymb,indentfirst,mathrsfs,amscd,url}
\usepackage[mathscr]{eucal} 
%\usepackage[active]{srcltx}




\setlength{\oddsidemargin}{0.4in}
\setlength{\evensidemargin}{0.4in}
\setlength{\textwidth}{5.5in}
\setlength{\textheight}{8.8in}
\setlength{\marginparwidth}{0.8in}
\addtolength{\headheight}{2.5pt}

\renewcommand{\thefootnote}{}

\begin{document}
\thispagestyle{empty}
\vspace*{-2cm}

 
\noindent
\begin{center}\Large\bfseries\textsf{\'Algebra 1 (grado en f\'isica)\\   curso 2013-14, primer semestre\\
	grupo 516}
\end{center}
\hrulefill{}
 
\noindent
{\sffamily    
Profesor:} Marco Zambon. \\
%{\sffamily Despacho:} 01.17.AU.103\\
{\sffamily Tutor\'ias:} por cita previa, en el despacho 605 del departamento de matem\'aticas (modulo 17). \\
{\sffamily P\'agina web del curso:} \url{http://www.uam.es/marco.zambon/WS13Phy.html}

\noindent\hrulefill{}  

\noindent
{\sffamily Horario:} Lunes a Jueves, 17:30-18:20. Aula  01.11.AU.201-4.   \\ 
\noindent
{\sffamily  Fecha de examenes:} \\ 
ejercicios realizados en clase: a decidir\\
%26 de octubre 2012 (viernes) y
% 30 de noviembre 2012 (viernes), en clase.\\ 
examen final: 8 de enero de 2014 (miercoles, tarde)\\
convocatoria extraordinaria:  20 de junio de 2014 (viernes, ma\~nana) 

\noindent\hrulefill{}  

\noindent
{\sffamily  Sistema de evaluaci\'on:}  (Ver la Guia Docente para m\'as detalles.)\\
La nota se determina a partir del siguiente promedio:\\
entrega de ejercicios realizados en clase (30 por ciento),\\
examen final (70 por ciento).

Los enunciados de los ejercicios realizados en clase consisteran de una selecci\'on de las hoja   para casa.

\noindent\hrulefill{}  

\noindent
{\sffamily  Hojas para casa:}\\
Habr\'a  hojas de ejercicios para casa, que discutiremos durante las horas de pr\'acticas. No ser\'a necesario entregar los ejercicios.
La participaci\'on durante las horas de pr\'acticas puede contribuir a subir la nota hasta un 10 por ciento.

%\noindent
%{\sffamily  Tareas para casa:} Cada 2 semanas habr\'a una hoja de ejercicios para casa, que discutiremos en clase en la semana siguiente. No ser\'a necesario entregar los ejercicios.

%\noindent\hrulefill{}  


\noindent\hrulefill{}  

 \noindent
{\sffamily Programa orientativo del curso:} \\
Ver la Guia Docente.
  

\noindent\hrulefill{}  

 \noindent
{\sffamily Bibliograf\'ia:} \\
Ver la Guia Docente para una lista completa de libros aconsejados. La referencia principal ser\'a:


\begin{itemize}
\item \emph{Sergei Treil. Linear algebra done wrong.} Disponible en la p\'agina web del curso. \end{itemize}

   
  

  
  
\end{document}
